\documentclass[11pt]{article}

\usepackage[margin=1in]{geometry}
\usepackage[T1]{fontenc}
\usepackage{lmodern}
\usepackage{microtype}
\usepackage{array}
\usepackage{longtable}
\usepackage{booktabs}
\usepackage{amsmath}
\usepackage{amsfonts}
\usepackage{xcolor}
\usepackage{url}

\definecolor{donecolor}{RGB}{20,105,62}
\definecolor{opencolor}{RGB}{160,55,35}
\definecolor{conditionalcolor}{RGB}{135,92,18}
\newcommand{\lean}[1]{\path{#1}}
\newcommand{\done}{\textcolor{donecolor}{locally complete}}
\newcommand{\openstate}[1]{\textcolor{opencolor}{open (#1)}}
\newcommand{\conditional}{\textcolor{conditionalcolor}{conditional}}

\makeatletter
\renewcommand\section{\@startsection{section}{1}{0pt}%
  {2.4ex plus .6ex minus .2ex}{1.0ex}%
  {\normalfont\Large}}
\renewcommand\subsection{\@startsection{subsection}{2}{0pt}%
  {1.8ex plus .4ex minus .2ex}{0.7ex}%
  {\normalfont\large}}
\makeatother

\setlength{\parindent}{0pt}
\setlength{\parskip}{0.55em}
\setlength{\emergencystretch}{3em}
\renewcommand{\arraystretch}{1.16}
\renewcommand{\descriptionlabel}[1]{\hspace{\labelsep}\normalfont #1}

\begin{document}

\begin{center}
  {\LARGE State and proof roadmap for \lean{Lemmas}}\\[0.6em]
  {\large Quantitative strict-AT formalization}\\[0.4em]
  Snapshot checked on 2 August 2026 at revision \lean{01cb6d5}
\end{center}

\section{Executive status}

The target \lean{LatalaMeetsAT} builds successfully with Lean 4.32.1.  The
latest check completed 3126 jobs.  The public theorem
\lean{SpinGlass.AT.quantitative_strictAT} has no analytic hypotheses in its
statement.  Its proof invokes the named existence theorems for the cavity
remainder and free-energy endpoints, both of which currently contain
\lean{sorry}.

The directory \lean{Lemmas} contains 24 Lean files.  Eighteen files have
no local \lean{sorry}; six files contain six occurrences in total.  There are
no project-specific \lean{axiom} declarations and no \lean{admit}s in this
directory.  Every placeholder is now in an ordinary theorem; there are no
placeholder definitions.

The former analytic typeclasses have been removed.  Gaussian covariance
differentiation, GT interpolation, the last-spin construction, the
cavity-remainder estimate, and the free-energy endpoint facts are isolated as
ordinary theorems with localized placeholders.  GT functional coercivity now
has a proved algebraic assembly theorem consuming the explicit
\lean{GTFunctionalAnalyticData} interface; its multiplier Taylor estimates and
negative-range gap still require analytic proofs.  The current
cavity-interpolation existence statement is discharged by a zero witness
because it does not encode the intended endpoint identities.
The propositions \lean{HasFreeEnergyEndpointBridge} and
\lean{HasCavityRemainderBound} remain useful interfaces for their proved
downstream consequences.

Explicit \lean{\#print axioms} checks show that both
\lean{SpinGlass.AT.quantitative_strictAT} and
\lean{SpinGlass.AT.SK.quantitative_strictAT_for_sk} depend on exactly
\lean{propext}, \lean{sorryAx}, \lean{Classical.choice}, and
\lean{Quot.sound}.  Thus the source is structurally build-complete but neither
public theorem is certified.

The wrapper \lean{Latala_AT.lean} no longer assumes
\lean{SmartPathBridge}.  Its endpoint overlap and free energy are defined
directly from \lean{RSSmartPathDisorder} at $s=1$, and all three conclusions
quantify directly over that path.  This is a bridge-free theorem at the
strict-AT model interface, not yet a conversion theorem for the separate
\lean{SpinGlass.SKDisorder} API.  That API has a different Hamiltonian sign
convention and does not itself supply the independent one-site field needed
to realize the complete path on the same probability space.

\section{Import spine}

The principal dependency spine is now

\begin{center}
\begin{minipage}{0.94\textwidth}\small
\lean{Replicas}
$\to$ \lean{RSParameters}
$\to$ \lean{UniformData}
$\to$ \lean{SmartPath}
$\to$ \lean{GaussianDerivative}
$\to$ \lean{FreeEnergyDerivative}
$\to$ \lean{Scalar/Semigroup}
$\to$ \lean{Scalar/LatalaKernel}
$\to$ \lean{Scalar/StrictATSign}
$\to$ \lean{GT/Defs}
$\to$ \lean{GT/Interpolation}
$\to$ \lean{GT/Coercivity}
$\to$ \lean{CoupledPressure}
$\to$ \lean{FixedDeviation}
$\to$ \lean{Absorption}
$\to$ \lean{FreeEnergy}
$\to$ \lean{Replicon}
$\to$ \lean{CompactCorollary}
$\to$ \lean{MainResult}.
\end{minipage}
\end{center}

The cavity branch starts independently at \lean{SmartPath}.  Definitions feed
both interpolation and coefficients; coefficients feed stability and the
system; absorption joins stability, the system, and fixed-deviation estimates.

\section{State of every Lean file}

\small
\begin{longtable}{@{}p{0.28\textwidth}p{0.15\textwidth}p{0.50\textwidth}@{}}
\toprule
File & Local state & What is present or missing \\
\midrule
\endhead
\lean{Replicas.lean} & \done & Finite configurations, positive partition functions, Gibbs-weight bounds and normalization, permutation invariance, overlap bounds, linearity, monotonicity, and the mixed-overlap estimate. \\
\lean{RSParameters.lean} & \done & Gaussian RS quantities, the selected fixed point, positivity and range of $q$, identities for $A$, and anomalous-versus-replicon comparison. \\
\lean{UniformData.lean} & \done & Compact-uniform parameter package and the strict-AT path gap. \\
\lean{SmartPath.lean} & \done & Covariance kernel, abstract centered Gaussian path, full Hamiltonian with external field, overlap second moment, and path gap.  It still has no explicit finite Gaussian coordinate realization. \\
\lean{GaussianDerivative.lean} & \openstate{1} & The covariance derivative and finite-replica contraction are explicit.  The finite-dimensional Gaussian interpolation theorem \lean{quenchedGibbs_deriv_of_covariance_deriv} remains. \\
\lean{FreeEnergyDerivative.lean} & \openstate{1} & Definitions and \lean{rsGap_deriv} are complete; \lean{smartPath_freeEnergy_deriv} remains. \\
\lean{Scalar/Semigroup.lean} & \done & Both branches of \lean{scalarPsi} and \lean{scalarTrialValue_eq} are proved, including Gaussian integrability and convolution. \\
\lean{Scalar/LatalaKernel.lean} & \done & Derivative and sign of \lean{latalaH}, antitonicity, the opposite-monotonicity covariance inequality, singular reference normalization, and the weighted kernel comparison are proved. \\
\lean{Scalar/StrictATSign.lean} & \conditional & \lean{scalarOrderParameterCorrect} is an explicit piecewise Gaussian formula.  The order-theoretic sign and local absolute-value estimate are proved from left and right linear bounds; the analytic derivation of those bounds is not supplied. \\
\lean{GT/Defs.lean} & \done & Attainable overlaps, constrained free energy, signed matrix path, mass profile, correction, terminal condition, and endpoint identities. \\
\lean{GT/Interpolation.lean} & \openstate{1} & Contains the zero, half, and one operators, the constant half-mass identity, explicit Gaussian recursion steps, \lean{gtSemigroupSolution}, and \lean{gtFunctional}.  The finite-volume theorem \lean{twoReplica_GT_bound} remains. \\
\lean{GT/Coercivity.lean} & \conditional & The local constant, abstract Taylor-loss lemma, overlap-range lemma, and algebraic coercivity assembly are proved.  \lean{GTFunctionalAnalyticData} records the remaining Taylor and negative-range analytic inputs. \\
\lean{CoupledPressure.lean} & \conditional & Correct quadratic pressure definitions are present.  The uniform sublinear overlap estimate is proved from \lean{fixedDeviation} by a supremum envelope and a small-tail argument. \\
\lean{FixedDeviation.lean} & \openstate{1} & Exponentially small fixed-deviation tails remain. \\
\lean{Cavity/Defs.lean} & \done & The observables $A,B,C$, the identity between $A$ and the overlap second moment, cavity vector, source vector, and $3\times3$ cavity matrix. \\
\lean{Cavity/Interpolation.lean} & \conditional & The third-moment bound is complete.  \lean{exists_cavityInterpolatedAverage} is proved by choosing the zero function, and the public definition and derivative estimate follow.  The statement lacks normalization and endpoint identities, so it does not construct the intended last-spin interpolation. \\
\lean{Cavity/Coefficients.lean} & \done & Replica-edge classification, decoupled coefficients, and equality with the displayed cavity matrix. \\
\lean{Cavity/System.lean} & \openstate{1} & The system is algebraic because the remainder is defined as the residual.  The uniform existence theorem \lean{exists_cavityRemainder_bound} remains. \\
\lean{Cavity/Stability.lean} & \done & Explicit matrix inverse with two-sided identities, triangular determinant formula, anomalous and replicon gaps, uniform inverse estimate, and replicon left eigenvector. \\
\lean{Absorption.lean} & \done & Pointwise cubic splitting, exact cavity-vector norm, domination of $B$ and $C$ by $A$, and the uniform $N A_s$ bound are proved from the stated cavity-remainder hypothesis. \\
\lean{FreeEnergy.lean} & \openstate{1} & The free-energy estimate is proved using the mean-value theorem, positivity, and absorption.  The uniform endpoint theorem \lean{freeEnergy_endpoint_bridge} remains. \\
\lean{Replicon.lean} & \done & Third-moment little-$o$ and replicon susceptibility are proved from earlier tail, absorption, stability, and assumed cavity-remainder results. \\
\lean{CompactCorollary.lean} & \done & Compactness plus continuity and strict AT produce \lean{UniformATData}. \\
\lean{MainResult.lean} & \conditional & The assembly has no analytic assumptions in its public statement and combines absorption, free energy, and replicon results through named existence theorems.  Those transitive placeholders still yield \lean{sorryAx}. \\
\bottomrule
\end{longtable}
\normalsize

\section{Remaining work and recommended prover}

The rating concerns the current source interfaces.  ``Codex'' should lead
definition design, API repair, supporting lemmas, and cross-file integration.
``Aristotle'' is useful for a focused Lean proof only after the relevant
interfaces and mathematical estimates have stable statements.

\small
\begin{longtable}{@{}p{0.25\textwidth}p{0.12\textwidth}p{0.13\textwidth}p{0.42\textwidth}@{}}
\toprule
File or contract & Difficulty & Better lead & Reason and possible handoff \\
\midrule
\endhead
Gaussian covariance derivative & difficult & Codex & The operator is written down, but \lean{quenchedGibbs_deriv_of_covariance_deriv} still requires finite-dimensional Gaussian covariance interpolation or an explicit affine realization with domination.  Hand off the final finite-sum algebra after the analytic API is fixed. \\
\lean{FreeEnergyDerivative.lean} & difficult & Codex & The proof must apply Gaussian covariance differentiation, derive the diagonal correction and covariance cancellation, and normalize replica sums.  Focused algebraic sublemmas may then be handed off. \\
\lean{Scalar/StrictATSign.lean} & difficult & Codex & The formula and order-theoretic closure are explicit.  The remaining analytic work is differentiation under Gaussian integrals and the left and right linear bounds. \\
GT interpolation theorem & difficult & Codex & The explicit recursion exists, but \lean{twoReplica_GT_bound} still requires the finite-cascade interpolation, endpoint identifications, and derivative sign. \\
GT functional analytic package & difficult & Codex & Construct \lean{GTFunctionalAnalyticData} by proving multiplier derivatives, a uniform Taylor bound, the scalar linear bounds, and the signed far-from-$q$ estimates.  The algebraic assembly \lean{gtFunctional_coercivity} is complete. \\
\lean{CoupledPressure.lean} & moderate & Codex & The sublinear overlap estimate is locally complete from \lean{fixedDeviation}.  Recheck its supremum construction after the concentration theorem is proved. \\
\lean{FixedDeviation.lean} & difficult & Codex & The covariance-only disorder still lacks the explicit coordinate or direct concentration interface needed for the gradient estimate. \\
\lean{Cavity/Interpolation.lean} & difficult & Codex & The last-spin decomposition is absent, and the derivative expansion requires replica bookkeeping, endpoint identities, and uniform estimates. \\
Free-energy endpoints & moderate & Codex & The final calculus argument is complete.  \lean{freeEnergy_endpoint_bridge} still needs continuity of the covariance-dependent free energy and the $s=0$ product endpoint identity. \\
Top-level certification & moderate & Codex & The assembly and bridge-free endpoint wrapper build.  Certification requires replacing the six theorem placeholders, repairing the cavity interpolation specification, and auditing the resulting axiom set. \\
\bottomrule
\end{longtable}
\normalsize

\section{Structure of the complete proof}

This section follows the mathematical order recorded in
\lean{blueprint_at.aux}.  Parenthetical references give the equation or
result number in the blueprint together with its auxiliary-file label.  The
proof has four principal stages:
\[
  \boxed{\text{RS coercivity}}
  \longrightarrow
  \boxed{\text{overlap concentration}}
  \longrightarrow
  \boxed{\text{cavity closure}}
  \longrightarrow
  \boxed{\text{quantitative conclusions}}.
\]
The logical role of each stage, and its present Lean status, is as follows.

\subsection{Parameters, fixed point, and smart path}

Fix a compact set $K$ of parameters $(\beta,h)$ in the strict AT region.  For
each $(\beta,h)\in K$, let $q=q(\beta,h)$ solve the replica-symmetric fixed
point equation (Blueprint (0.1), \lean{eq:q}), and set
\[
  \alpha(\beta,h)
  =\beta^2\,\mathbb E\!\left[\operatorname{sech}^4
    (h+\beta\sqrt q\,Z)\right].
\]
Strict AT means that $\alpha<1$; compactness gives a uniform gap
$1-\alpha\geq\delta>0$ (Blueprint (0.3)--(0.6),
\lean{eq:alpha}--\lean{eq:qcompact}).  Interpolate between the independent
one-site random-field model at $s=0$ and the SK Hamiltonian at $s=1$ by the
replica-symmetric smart path (Blueprint (0.7), \lean{eq:path}).  Write
\[
  A_s=\mathbb E\langle(R_{12}-q)^2\rangle_s,
  \qquad
  D_N(s)=p_{\mathrm{RS}}(s)-p_N(s).
\]
The finite spin space, Gibbs measure, fixed-point identities, compact
uniformity package, and covariance kernel of this path are formalized in
\lean{Replicas.lean}, \lean{RSParameters.lean}, \lean{UniformData.lean}, and
\lean{SmartPath.lean}.

The first analytic identity is Gaussian covariance differentiation.  It
gives
\[
  D_N'(s)=\frac{\beta^2}{4}A_s,
  \qquad 0<s<1,
\]
which is Blueprint (2.7), \lean{eq:DNprime}, and is also the derivative form
of the free-energy identity used at the end.  The general Gaussian
integration-by-parts rule (Blueprint (1.1), \lean{eq:GIBP}) and this
specialized derivative identity are not yet verified in Lean; they are the
placeholders in \lean{GaussianDerivative.lean} and
\lean{FreeEnergyDerivative.lean}.

\subsection{Stage I: replica-symmetric coercivity}

The purpose of this stage is to prove that forcing the overlap away from $q$
costs a quantity of order $(v-q)^2$, uniformly on $K$ and along the whole
smart path.

First solve the scalar replica-symmetric Parisi equation by its explicit
Gaussian semigroup (Blueprint (1.4)--(1.9),
\lean{eq:parisipde}--\lean{eq:RSpathvalue}).  The strict AT condition must then
be converted into the one-sided sign estimate for the scalar order parameter
$g_s$:
\[
  g_s(u)-u\geq\delta(q-u)\quad(0\leq u\leq q),
  \qquad
  u-g_s(u)\geq\delta(u-q)\quad(q\leq u\leq1).
\]
These are the linear form of Blueprint (1.13)--(1.18), in particular
\lean{eq:linearATsign}, \lean{eq:leftstrict}, and \lean{eq:fprime}.  Lean
contains the explicit scalar formula and proves all order-theoretic
consequences of these two inequalities, but the analytic derivation of the
inequalities from strict AT is still missing.

Next constrain two replicas to have overlap $v$ and introduce the
two-dimensional Guerra--Talagrand functional
$\mathcal P_{\mathrm{GT}}(\lambda,v)$ (Blueprint (1.19)--(1.26),
\lean{eq:constrainedZ}--\lean{eq:GTfunctional}).  The specialized
Guerra--Talagrand interpolation gives
\[
  \mathbb E p_{N,s}^{(2)}(v)
  \leq \mathcal P_{\mathrm{GT}}(\lambda,v)
\]
(Blueprint (1.28) and Appendix A, \lean{eq:GTfunctionalbound} and
\lean{eq:specialGTlemma}).  Taylor expansion in $\lambda$, the scalar sign
estimate near $v=q$, and a separate signed-path estimate for negative
overlaps then yield
\[
  \mathbb E p_{N,s}^{(2)}(v)
  \leq 2p_{\mathrm{RS}}(s)-c(v-q)^2
\]
(Blueprint (1.29), \lean{eq:GTcoercivity}).  The explicit recursion and the
algebraic Taylor minimization are formalized.  The finite-cascade GT bound,
the multiplier Taylor estimates, and the uniform negative-overlap gap
(Blueprint (1.36)--(1.52)) remain unverified.

\subsection{Stage II: preliminary overlap concentration}

Introduce a small quadratic coupling of two replicas and its normalized
coupled pressure (Blueprint (2.1)--(2.3), \lean{eq:lambda0},
\lean{eq:quadraticpressure}, and \lean{eq:normalizedcoupling}).  The coercive
bound from Stage I controls the pressure away from $q$.  Combining that bound
with the derivative identity for $D_N$, and closing the resulting integral
inequality by Gronwall's lemma, gives the preliminary estimate in Blueprint
(2.8)--(2.9), \lean{eq:gronwallDN} and \lean{eq:preconcentration}.

Gaussian concentration of the constrained log-partition functions then
upgrades this control to the fixed-deviation tail estimate: for every
$\varepsilon>0$ there exist $c,C>0$ such that
\[
  \mathbb E\Bigl\langle
    \mathbf 1_{\{|R_{12}-q|\geq\varepsilon\}}
  \Bigr\rangle_s
  \leq Ce^{-cN}
\]
uniformly in $(\beta,h)\in K$ and $s\in[0,1]$ (Blueprint (2.10),
\lean{eq:tail}).  The deterministic pressure deduction is formalized in
\lean{CoupledPressure.lean}, but its present Lean proof assumes the
fixed-deviation estimate and therefore runs in the reverse direction from the
blueprint.  The direct pressure-envelope proof and the Gaussian concentration
step are not yet formalized; the latter is the placeholder in
\lean{FixedDeviation.lean}.

\subsection{Stage III: cavity closure}

Remove the last spin and interpolate between the $N$-spin model and its
$(N-1)$-spin cavity decomposition.  For the three overlap observables
\[
  \mathbf v_s=(A_s,B_s,C_s)^{\mathsf T},
\]
the cavity expansion gives the finite-dimensional system
\[
  (I-sL)\mathbf v_s=\frac1N\,\boldsymbol\theta+\boldsymbol\rho_s
\]
(Blueprint (3.1)--(3.5), \lean{eq:cavitysystem}--\lean{eq:rho}).  The
last-spin interpolation, replica differentiation, endpoint factorization,
and Taylor comparison (Blueprint (3.6)--(3.22)) must prove the uniform
remainder estimate
\[
  \|\boldsymbol\rho_s\|
  \leq C\left(N^{-3/2}+
    \mathbb E\langle|R_{12}-q|^3\rangle_s\right).
\]
Lean verifies the cavity observables, coefficients, and displayed linear
identity, but the present interpolation is under-specified and the remainder
bound remains a placeholder in \lean{Cavity/System.lean}.

The matrix $I-sL$ is invertible throughout the strict AT region.  Its
triangularization separates the anomalous and replicon factors, and the
uniform AT gap bounds its inverse (Blueprint (3.23)--(3.26),
\lean{eq:gammas}--\lean{eq:inverse}).  This matrix algebra, including the
replicon left eigenvector, is fully verified in \lean{Cavity/Stability.lean}.

\subsection{Stage IV: absorption and the three conclusions}

Apply the inverse bound to the cavity system.  Split the cubic moment at a
small fixed threshold and use the exponential tail from Stage II.  The cubic
term can then be absorbed into the left side, giving
\[
  \sup_{(\beta,h)\in K}\sup_{s\in[0,1]} N A_s<\infty
\]
(Blueprint (4.1)--(4.2), \lean{eq:preabsorb} and \lean{eq:thirdsplit}).  This
absorption argument is verified in \lean{Absorption.lean}, conditional only
on the unproved fixed-deviation and cavity-remainder estimates.

The uniform $A_s=O(N^{-1})$ estimate has two consequences.  First, integrate
$D_N'(s)=\beta^2A_s/4$ from $0$ to $1$.  Using continuity and the product
identity $D_N(0)=0$ gives
\[
  0\leq p_{\mathrm{RS}}-p_N\leq \frac{C}{N}
\]
(Blueprint (4.3), \lean{eq:freeenergyidentity}).  The integration argument is
formalized, but continuity and the $s=0$ endpoint identification remain the
placeholder in \lean{FreeEnergy.lean}.

Second, fixed-deviation concentration and $N A_s=O(1)$ imply
\[
  N\,\mathbb E\langle|R_{12}-q|^3\rangle_s\longrightarrow0
\]
uniformly.  Apply the replicon vector $(1,-2,1)$ to the cavity system.  Since
it is a left eigenvector of $L$, one obtains
\[
  N(A_s-2B_s+C_s)
  \longrightarrow
  \frac{a(\beta,h)}{1-s\alpha(\beta,h)}
\]
uniformly (Blueprint (4.4)--(4.5), \lean{eq:repliconeq} and
\lean{eq:repliconrho}).  This final linear-algebraic deduction is verified in
\lean{Replicon.lean}, again conditional on the earlier tail and remainder
inputs.

\subsection{Final assembly and certification boundary}

Compactness converts pointwise strict AT into the uniform data used above.
The theorem \lean{SpinGlass.AT.quantitative_strictAT} then packages exactly
the three conclusions: $N A_s=O(1)$, an $O(N^{-1})$ replica-symmetric
free-energy error, and the uniform replicon susceptibility limit.  This
assembly is verified in \lean{MainResult.lean}.  Its dependency chain is
nevertheless uncertified because it invokes the unproved cavity-remainder
and endpoint theorems, which in turn depend on the missing Gaussian, GT,
concentration, and last-spin arguments.  In the terminology of Appendix B of
the blueprint, certification requires the audit in subsection B.36 and the
resolution of the high-risk obligations in subsection B.37, as recorded in
\lean{blueprint_at.aux}.

\section{Mathematical statements not yet verified in Lean}

This section states the missing results in ordinary mathematical language.
Write $R_{12}$ for the overlap of two Gibbs replicas, $q$ for the
replica-symmetric fixed point, and
\[
  A_s=\mathbb E\bigl\langle (R_{12}-q)^2\bigr\rangle_s,
  \qquad
  T_s=\mathbb E\bigl\langle |R_{12}-q|^3\bigr\rangle_s.
\]
The expectation includes both the Gaussian disorder and the Gibbs measure.
The constants below are required to be uniform when $(\beta,h)$ ranges over
the fixed compact strict-AT parameter set and $s\in[0,1]$.

\subsection{The six assertions currently containing \lean{sorry}}

\begin{enumerate}
\item \lean{GaussianDerivative.lean:quenchedGibbs_deriv_of_covariance_deriv}.
Let $(H_s(\sigma))_{\sigma\in\{-1,1\}^N}$ be a centered finite-dimensional
Gaussian process whose covariance matrix $C_s(\sigma,\tau)$ is affine in
$s$.  For every bounded observable $F$ of $n$ replicas, prove that
$s\mapsto\mathbb E\langle F\rangle_s$ is differentiable for $0<s<1$ and
that its derivative is one half of the contraction of
$\dot C_s$ with the Hessian of the normalized Gibbs expectation.  In replica
notation this is the standard formula involving the original $n$ replicas
and two additional replicas, including all terms arising from
differentiating the partition function.  The Lean file defines this
contraction explicitly as \lean{replicaCovarianceOperator}; what is missing
is the Gaussian integration-by-parts argument proving the formula.

\item \lean{FreeEnergyDerivative.lean:smartPath_freeEnergy_deriv}.
For
\[
  p_N(s)=\frac1N\,\mathbb E\log Z_N(s),
\]
prove, for $0<s<1$, the smart-path identity
\[
  p_N'(s)=\frac{\beta^2}{4}
  \left((1-q)^2-
  \mathbb E\bigl\langle(R_{12}-q)^2\bigr\rangle_s\right).
\]
This requires applying the preceding Gaussian interpolation formula and
checking the diagonal term, the covariance cancellation, and the
normalization of the replica sums.

\item \lean{GT/Interpolation.lean:twoReplica_GT_bound}.
For every attainable overlap $v$, multiplier $\lambda\in\mathbb R$, and
$s\in[0,1]$, prove the two-replica Guerra--Talagrand upper bound
\[
  \mathbb E\,p_{N,s}^{(2)}(v)
  \leq \mathcal P_{\mathrm{GT}}(\beta,h,q,s,\lambda,v),
\]
where the left side is the free energy of two replicas constrained to have
overlap $v$, and the right side is the explicit finite Gaussian recursion
defined in \lean{GT/Interpolation.lean}.  The missing mathematics is the
finite-cascade interpolation, including identification of both endpoints
and proof that its derivative has the required sign.

\item \lean{FixedDeviation.lean:fixedDeviation}.
For every $\varepsilon>0$, prove that there are constants $c,C>0$, independent
of $N$, $s$, and $(\beta,h)$ in the chosen compact set, such that
\[
  \mathbb E\Bigl\langle
    \mathbf 1_{\{|R_{12}-q|\geq\varepsilon\}}
  \Bigr\rangle_s
  \leq C e^{-cN}.
\]
The intended proof combines the quadratic loss for constrained overlap with
a Gaussian concentration inequality.  In particular it still needs the
uniform $O(N)$ bound for the squared Euclidean gradient of the relevant
log-partition functions in Gaussian coordinates.

\item \lean{Cavity/System.lean:exists_cavityRemainder_bound}.
Let $\mathbf r_s$ denote the error in the three-component cavity system.
Prove that there is a constant $C>0$, uniform in the above parameters, for
which
\[
  \|\mathbf r_s\|
  \leq C\left(N^{-3/2}+T_s\right).
\]
The displayed linear cavity system is already an algebraic identity because
$\mathbf r_s$ is defined as its residual.  The unverified content is this
quantitative bound, which should follow from a genuine last-spin
interpolation and the proved coefficient identities.

\item \lean{FreeEnergy.lean:freeEnergy_endpoint_bridge}.
Define
\[
  D_N(s)=p_{\mathrm{RS}}(s)-p_N(s),
\]
where $p_{\mathrm{RS}}(s)$ is the replica-symmetric trial value along the
smart path.  Prove that $D_N$ is continuous on $[0,1]$ and that
\[
  D_N(0)=0.
\]
Equivalently, identify the $s=0$ Gaussian process with the independent
one-site random-field model and factorize its partition function.  This is
the endpoint input needed to integrate the already stated derivative
identity up to $s=1$.
\end{enumerate}

\subsection{Further analytic content not supplied by the current files}

The absence of \lean{sorry} in the following files does not mean that all of
the intended mathematics has been established.  In these cases the missing
content is hidden in an explicit hypothesis or is absent from the theorem
statement.

\begin{enumerate}
\item For the scalar order-parameter map $g_s(u)$, prove the uniform one-sided
linear estimates
\[
  g_s(u)-u\geq \delta(q-u)\quad(0\leq u\leq q),
  \qquad
  u-g_s(u)\geq \delta(u-q)\quad(q\leq u\leq1),
\]
for some $\delta>0$.  Lean verifies that these inequalities imply the desired
sign of $g_s(u)-u$ and its local linear separation from zero, but it does not
derive the inequalities from the strict AT condition.

\item For the Guerra--Talagrand functional, prove the analytic estimates used
in its coercivity argument.  On $-q\leq v\leq1$, this means a uniform
second-order upper Taylor bound in the multiplier $\lambda$, together with a
lower bound proportional to $|v-q|$ for the magnitude of the first
derivative at $\lambda=0$.  On $-1\leq v\leq-q$, prove a uniform strictly
positive gap below twice the replica-symmetric trial value.  Lean proves the
quadratic coercivity conclusion from these assumptions, but no theorem
constructs the assumptions from the explicit recursion.

\item Construct the actual last-spin interpolation $\nu_{s,u}(F)$, including
its normalization and its identities at $u=0$ and $u=1$, and prove
\[
  \left|\frac{d^2}{du^2}\nu_{s,u}(F)\right|
  \leq C\left(T_s+N^{-3}\right)
  \qquad (|F|\leq1).
\]
The present existence theorem is formally true only because it chooses the
identically zero function; its statement imposes no endpoint conditions and
therefore does not yet describe the cavity interpolation used in the
mathematical argument.
\end{enumerate}

Consequently, the two public strict-AT theorems are not Lean-certified even
though their final assembly proofs contain no local placeholders: they depend
transitively on the six unproved assertions above.  Completing only those six
proofs would remove the current occurrences of \lean{sorry}, but a faithful
formalization of the intended proof also requires the three additional
analytic items in this subsection.

\section{Named analytic interfaces}

\begin{description}
\item[Scalar linear bounds.]
\lean{strictAT_sign_of_scalar_linear_bounds} proves the required sign and
local absolute-value conclusion once the left and right analytic inequalities
for \lean{scalarOrderParameterCorrect} are supplied.  No theorem currently
derives those inequalities from \lean{UniformATData}.

\item[Free-energy endpoints.]
\lean{HasFreeEnergyEndpointBridge path} packages continuity on $[0,1]$ and the
$s=0$ product identity.  The theorem \lean{freeEnergy_endpoint_bridge} is the
localized placeholder that supplies it uniformly.

\item[Cavity remainder.]
\lean{HasCavityRemainderBound data C} is the analytic input to absorption and
the final conclusions.  The theorem \lean{exists_cavityRemainder_bound} is the
localized placeholder that supplies such a constant.

\item[GT functional coercivity.]
\lean{GTFunctionalAnalyticData} packages the multiplier Taylor estimate and
the uniform negative-range gap required by the proved algebraic assembly
\lean{gtFunctional_coercivity}.  The package itself still needs analytic
construction from the explicit GT recursion.
\end{description}

These interfaces make dependencies visible.  Their placeholder existence
theorems and the GT analytic package must be supplied before certification.

\section{Recommended proof roadmap}

\begin{description}
\item[Model and Gaussian API.] Prove the mathematical claim represented by
\lean{quenchedGibbs_deriv_of_covariance_deriv}: if $H_s$ is a centered finite
Gaussian process with differentiable covariance $C_s$, then for every
observable $F$ of $n$ replicas,
\[
 \frac d{ds}\mathbb E\langle F\rangle_s
 =\frac12\mathbb E\left\langle F\left(
 \sum_{a,b\le n}\dot C_s(\sigma^a,\sigma^b)
 -2n\sum_{a\le n}\dot C_s(\sigma^a,\sigma^{n+1})
 +n(n+1)\dot C_s(\sigma^{n+1},\sigma^{n+2})
 -n\dot C_s(\sigma^{n+1},\sigma^{n+1})\right)\right\rangle_s.
\]
Construct this result from a finite-coordinate Gaussian-law module, or give
the smart path an explicit affine realization.  Specializing to $F=1$ after
differentiating the logarithmic partition function must then give the claim
in \lean{smartPath_freeEnergy_deriv},
\[
 p_N'(s)=\frac{\beta^2}{4}\bigl((1-q)^2-A_s\bigr),
 \qquad
 D_N'(s)=\frac{\beta^2}{4}A_s.
\]
Use the same Gaussian coordinates for concentration and the last-spin
decomposition.

\item[Scalar package.] Establish the two analytic claims assumed by
\lean{strictAT_sign_of_scalar_linear_bounds}:
\[
 g_s(u)-u\geq\delta(q-u)\quad(0\leq u\leq q),
 \qquad
 u-g_s(u)\geq\delta(u-q)\quad(q\leq u\leq1).
\]
The theorem bearing that name proves from these estimates that $g_s(u)-u$ is
positive to the left of $q$, negative to its right, and satisfies
$|g_s(u)-u|\geq c|u-q|$ locally.  The order-theoretic closure, convolution
identity, and Lata\l a kernel lemmas are complete; the Gaussian analytic
derivation of the two displayed inequalities is not.

\item[GT package.] Prove the claim \lean{twoReplica_GT_bound}:
\[
 \mathbb E p_{N,s}^{(2)}(v)
 \leq \mathcal P_{\mathrm{GT}}(\beta,h,q,s,\lambda,v)
\]
for every attainable overlap $v$ and multiplier $\lambda$.  Next construct
\lean{GTFunctionalAnalyticData}, whose mathematical content is a uniform
Taylor upper bound
\[
 \mathcal P_{\mathrm{GT}}(\lambda,v)
 \leq \mathcal P_{\mathrm{GT}}(0,v)+d(v)\lambda
       +\frac M2\lambda^2,
 \qquad |d(v)|\geq\delta|v-q|,
\]
for $-q\leq v\leq1$, together with a constant $\kappa>0$ such that for
$-1\leq v\leq-q$ some multiplier satisfies
\[
 \mathcal P_{\mathrm{GT}}(\lambda,v)
 \leq2p_{\mathrm{RS}}(s)-\kappa.
\]
The proved theorem \lean{gtFunctional_coercivity} then yields the precise
claim
\[
 \exists c>0\ \forall v\in[-1,1]\ \exists\lambda:\quad
 \mathcal P_{\mathrm{GT}}(\lambda,v)
 \leq2p_{\mathrm{RS}}(s)-c(v-q)^2.
\]

\item[Pressure and tails.] Prove directly the blueprint pressure-envelope
claim
\[
 p_{N,s}^{(2)}(\lambda_0)\leq p_{\mathrm{RS}}(s)+\varepsilon_N,
 \qquad \varepsilon_N\longrightarrow0,
\]
using GT coercivity, the fact that there are at most $N+1$ attainable
overlaps, and Gaussian concentration.  Combine this with
$D_N'=\beta^2A_s/4$ and Gronwall's inequality to obtain
$\sup_s A_s=o(1)$.  Finally prove \lean{fixedDeviation}, namely that for each
$\varepsilon>0$ there exist $c,C>0$ for which
\[
 \mathbb E\langle\mathbf1_{\{|R_{12}-q|\geq\varepsilon\}}\rangle_s
 \leq Ce^{-cN}.
\]
The current \lean{CoupledPressure.lean} contains the finite overlap
decomposition and cardinality bound, but it does not contain a theorem named
\lean{coupledPressure_sublinear}; the roadmap should not cite that obsolete
name.

\item[Cavity package.] Define the genuine last-spin interpolation
$\nu_{s,u}(F)$ with normalization and endpoint identities, and prove
\[
 \left|\frac{d^2}{du^2}\nu_{s,u}(F)\right|
 \leq C\left(T_s+N^{-3}\right),
 \qquad |F|\leq1.
\]
Use it to establish the mathematical claim in
\lean{exists_cavityRemainder_bound}: for the residual $\boldsymbol\rho_s$ in
the three-dimensional cavity system,
\[
 \|\boldsymbol\rho_s\|
 \leq C\left(N^{-3/2}+T_s\right)
\]
uniformly on the compact strict-AT set.  This is the missing analytic link
between the cavity expansion and the verified absorption argument.

\item[Free-energy endpoints.] Prove the claim
\lean{freeEnergy_endpoint_bridge}: the function
\[
 D_N(s)=p_{\mathrm{RS}}(s)-p_N(s)
\]
is continuous on $[0,1]$ and satisfies $D_N(0)=0$.  The proof should identify
the $s=0$ Gaussian law with independent one-site fields and factorize its
partition function.  Once this claim and the remainder estimate hold, the
verified calculus gives
\[
 0\leq p_{\mathrm{RS}}-p_N\leq C/N.
\]

\item[Concrete SK API.] If a theorem for an arbitrary
\lean{SpinGlass.SKDisorder} is required, prove a law-preserving conversion to
a product-space smart path, including an independent one-site Gaussian field
and reconciliation of the Hamiltonian signs.  The claim currently proved by
\lean{quantitative_strictAT_for_sk} concerns an
\lean{RSSmartPathDisorder} and asserts, at its $s=1$ endpoint,
\[
 NA_1\leq C,
 \qquad 0\leq p_{\mathrm{RS}}-p_N\leq C/N,
\]
together with the uniform replicon limit
\[
 N(A_s-2B_s+C_s)\longrightarrow
 \frac{a(\beta,h)}{1-s\alpha(\beta,h)}.
\]
It is not yet a conversion theorem for an independently supplied concrete SK
disorder.

\item[Certification.] Run \lean{lake build LatalaMeetsAT}, search the Lean
sources for \lean{sorry}, \lean{admit}, and project-specific \lean{axiom}, and
inspect \lean{\#print axioms SpinGlass.AT.quantitative_strictAT} and
\lean{\#print axioms SpinGlass.AT.SK.quantitative_strictAT_for_sk}.
Certification requires a successful build and absence of \lean{sorryAx}; it
also requires that the public theorems no longer obtain their analytic content
from unproved assumptions.
\end{description}

\section{Structural issues}

\begin{itemize}
\item[Cavity system.] The theorem \lean{cavity_system} follows by unfolding the
definition of \lean{cavityRemainder}.  Its analytic content is the separate,
unconstructed remainder bound.

\item[Cavity interpolation.] The theorem
\lean{exists_cavityInterpolatedAverage} currently chooses the constant zero
function.  This satisfies the displayed derivative bound because the theorem
does not require normalization or the $u=0,1$ endpoint identities.  The
specification must be strengthened before it represents the paper's
last-spin interpolation.

\item In \lean{GaussianDerivative.lean}, the theorem
\lean{quenchedGibbs_deriv_of_covariance_deriv} is the localized placeholder.
It does not yet prove Gaussian interpolation from the covariance assumptions
in \lean{SmartPath.lean}.

\item \lean{GT/Interpolation.lean:twoReplica_GT_bound} remains a localized
placeholder.  In \lean{GT/Coercivity.lean}, the explicit recursion and the
composition theorem \lean{gt_quadratic_coercivity} are present, and the
coercivity assembly is proved under \lean{GTFunctionalAnalyticData}; the
analytic package itself is not yet constructed.

\item \lean{FreeEnergy.lean} contains the placeholder theorem
\lean{freeEnergy_endpoint_bridge}.  The subsequent mean-value argument and
public free-energy estimate are complete from that result.

\item \lean{MainResult.lean:quantitative_strictAT} exposes no endpoint or
cavity hypothesis.  Its proof calls \lean{freeEnergy_endpoint_bridge} and
\lean{exists_cavityRemainder_bound}, so it remains uncertified until those and
the other transitive placeholders are proved.

\item \lean{Latala_AT.lean:quantitative_strictAT_for_sk} has no
\lean{SmartPathBridge} argument.  It quantifies directly over
\lean{RSSmartPathDisorder} and reads the SK quantities at $s=1$.  It should not
be described as a theorem for an arbitrary concrete
\lean{SpinGlass.SKDisorder} until a law-preserving model conversion is proved.

\item \lean{SmartPath.lean} specifies Gaussian marginals and covariance but no
concrete coordinate realization.  This single model-level gap blocks the
Gaussian derivative theorem, direct concentration, endpoint-law continuity,
and the last-spin decomposition.
\end{itemize}

\end{document}
